\clearpage\fancyhead[L]{\small PALLAS.}\SourceLeaf{381}
\begin{center}\textbf{PALLAS.}\par\vspace{4pt}\hrule\vspace{4pt}\fontsize{10}{12}\selectfont Astronomisches Jahrbuch für 1817. pp. 212\;..\;215. Berlin 1814.\par\vspace{4pt}\hrule\vspace{4pt}\end{center}

\begin{flushright}
Göttingen, 29 May 1814.
\end{flushright}

\par\indent My most esteemed friend, I still owe you my warmest thanks for your valuable gift of the astronomical yearbook for 1816. I delayed my thanks, hoping to send them along with some contributions to the new yearbook. The work I wanted to send you on the last opposition of Vesta has been held up, but I will send it to you later.

\Needspace{14\baselineskip}\par\indent I have already told you something about the last opposition of Juno. Since I can now add several things and no longer remember exactly how much I have already written to you, I will bring together everything essential here. The weather allowed me to make only one observation around the time of opposition, though I consider it a very good one. It was nearly three weeks before I could make a second observation as a check, and I include this second observation as reduced by Mr. \textsc{Nicolai}:

\begin{center}\fontsize{10}{12}\selectfont\setlength{\tabcolsep}{3pt}\renewcommand{\arraystretch}{1.15}\begin{tabular}{l|c|rrr|rrr}
1813 & Mean time & \multicolumn{3}{c|}{Apparent R.A.} & \multicolumn{3}{c}{Southern decl.}\\\hline
Nov. 19 & $9^{\mathrm u}\;46^{\mathrm m}\;46^{\mathrm s}$ & $60^\circ$ & $30'$ & $35''0$ & $3^\circ$ & $5'$ & $9''3$\\
Dec. 9 & $7\;9\;38$ & 56 & 56 & 47.4 & 3 & 56 & 35.3\\
\end{tabular}\end{center}

\Needspace{14\baselineskip}\par\indent I had already based a preliminary calculation of the opposition on my first observation. But since I had only calculated it hastily and \textsc{Bessel} subsequently sent me his meridian observations (which you will no doubt have received from him directly), I asked a skilful and talented student of mine, Mr. \textsc{Möbius}, to revise the results accordingly. He finds the opposition as follows:

\begin{center}19 November 1813. $18^{\mathrm u}\;12^{\mathrm m}\;0^{\mathrm s}$ Mean time in Göttingen\end{center}
\begin{center}\fontsize{10}{12}\selectfont\setlength{\tabcolsep}{2pt}\renewcommand{\arraystretch}{1.15}\begin{tabular}{lcr}
True longitude & $.\;.\;.\;.\;.\;.\;.\;.\;.\;.\;.\;.\;.\;.$ & $57^\circ\;34'\;2''1$\\
Geocentric latitude & $.\;.\;.\;.\;.\;.\;.\;.\;.\;.\;.$ & $23\;18\;46.1\;\mathrm{S.}$\\
\end{tabular}\end{center}

\Needspace{14\baselineskip}\par\indent At the same time he fitted the elements to this opposition and the three preceding ones of 1810, 1811 and 1812, and obtained the following result:

\begin{center}\fontsize{9}{11}\selectfont\setlength{\tabcolsep}{2pt}\renewcommand{\arraystretch}{1.15}\begin{tabular}{lcr}
Epoch 1810, Göttingen meridian & $.\;.\;.\;.\;.\;.\;.$ & $95^\circ\;29'\;53''2$\\
Daily mean tropical motion & $.\;.\;.\;.\;.\;.\;.$ & $812''7140$\\
Longitude of perihelion (1810 sid. ruh.) & $.\;.\;.\;.\;.\;.\;.$ & $53^\circ\;6'\;43''0$\\
Longitude of ascending node (1810) & $.\;.\;.\;.\;.\;.$ & $171\;6\;45.0$\\
Eccentricity angle & $.\;.\;.\;.\;.\;.\;.\;.\;.\;.\;.\;.\;.$ & $14\;43\;9.5$\\
Inclination of orbit & $.\;.\;.\;.\;.\;.\;.\;.\;.\;.\;.\;.\;.\;.$ & $13\;4\;12.9$\\
Logarithm of semimajor axis & $.\;.\;.\;.\;.\;.\;.$ & $0.4267616$\\
\end{tabular}\end{center}

\par\indent This skilful calculator has also computed an ephemeris for the next apparition of Juno, and I have the pleasure of enclosing a copy.

\SourceLeaf{382}\par\indent For Pallas, Mr. \textsc{Nicolai} (who has been employed for half a year as an assistant at the Seeberg Observatory) has likewise computed an ephemeris and sent me the enclosed copy for your yearbook. It is certainly more precise than necessary for finding the planet (which will in any case be fairly bright this time), even though it was computed from the earlier elements without taking account of the changes in them caused by Jupiter’s continuing perturbations. It is interesting all the same to state \emph{in advance} what effect these perturbations will have, so I assigned this rather laborious task to Mr. \textsc{Möbius}. The result is that around the time of opposition the right ascension in Mr. \textsc{Nicolai}’s ephemeris must be reduced by 1$\frac{1}{2}$ minutes, while the declination decreases by only a few seconds. For the opposition itself he obtains the following prediction:

\begin{center}25 October 1814. $12^{\mathrm u}\;39^{\mathrm m}\;50^{\mathrm s}$ Göttingen time.\end{center}
\begin{center}\fontsize{10}{12}\selectfont\setlength{\tabcolsep}{2pt}\renewcommand{\arraystretch}{1.15}\begin{tabular}{lcr}
True longitude & $.\;.\;.\;.\;.\;.\;.\;.\;.\;.$ & $31^\circ\;58'\;28''4$\\
Heliocentric latitude & $.\;.\;.\;.\;.\;.\;.$ & $23\;38\;32.7\;\mathrm{S.}$\\
\end{tabular}\end{center}

\noindent (He forgot to send me the geocentric latitude. As he is no longer here, I cannot include it for comparison with Mr. \textsc{Nicolai}’s result.)

\par\indent I have now begun to calculate the perturbations of Pallas due to Mars as well. Almost all will be very small, but they will be numerous. The correction terms due to Jupiter, Saturn and Mars that exceed $0''1$ will number more than 1000.

\Needspace{10\baselineskip}\par\indent A couple of days ago I received an excellent heliometer from Munich, with a focal length of 43 inches and an aperture of 34 lines. The ingenious instrument maker made several improvements, including the important step of adapting it for repeated measurements by making the two halves of the objective independently movable. This entirely avoids the detrimental effect of small local irregularities in the screw. The telescope’s precision is exceptionally high, as one has come to expect from \textsc{Frauenhofer}’s work.
